\subsection{Case of a closed central subgroup}

In this issue, we assume that the group G is unimodular.

 Consider a closed subgroup Z of the center of G, and let dz denote a Haar measure on Z. We endow the quotient group G/Z with the Haar measure \(\nu = \mu/dz\) . Let \(\chi\) a unitary character of Z. The quotient group G/Z is unimodular by INT, VII, p. 61, § 2, no. 7, cor. of prop. 11.

Lemma 7. --- One has

\[
\mathrm{N} _ {1} (\check {f}) = \mathrm{N} _ {1} (f), \quad \langle f _ {1} | f _ {2} \rangle = \langle \check {f} _ {1} | \check {f} _ {2} \rangle .\tag{6}
\]

for all \(f \in \mathcal{F}_{\chi}(\mathrm{G})\) and for all \(f_1\) and \(f_2\) in \(\mathcal{L}_{\chi}^{2}(\mathrm{G})\) .

These formulas are consequences of the definitions.

 For every \(g \in \mathbf{G}\) and every \(f \in \mathcal{F}_{\chi}(\mathbf{G})\) , the functions \(x \mapsto f(g^{-1}x)\) and \(x \mapsto f(xg)\) belong to \(\mathcal{F}_{\chi}(\mathbf{G})\) . They are denoted respectively \(\boldsymbol{\gamma}_{\mathrm{G},\chi}(g)f\) and \(\boldsymbol{\delta}_{\mathrm{G},\chi}(g)f\) . The maps \(\boldsymbol{\gamma}_{\mathrm{G},\chi}\) and \(\boldsymbol{\delta}_{\mathrm{G},\chi}\) are linear representations of \(\mathbf{G}\) in \(\mathcal{F}_{\chi}(\mathbf{G})\) . For every \(z \in \mathbf{Z}\) , we have

\[
\boldsymbol {\gamma} _ {\mathrm{G}, \chi} (g z) = \overline {{\chi (z)}} \boldsymbol {\gamma} _ {\mathrm{G}, \chi} (g), \quad \boldsymbol {\delta} _ {\mathrm{G}, \chi} (g z) = \chi (z) \boldsymbol {\delta} _ {\mathrm{G}, \chi} (g).\tag{7}
\]

Let \(f\in \mathcal{F}_{\chi}(\mathrm{G})\) and \(g\in \mathrm{G}\) ; we have

\[
| \boldsymbol {\gamma} _ {\mathrm{G}, \chi} (g) f | = \boldsymbol {\gamma} _ {\mathrm{G} / \mathrm{Z}} (g \mathrm{Z}) | f |, \quad | \boldsymbol {\delta} _ {\mathrm{G}, \chi} (g) f | = \boldsymbol {\delta} _ {\mathrm{G} / \mathrm{Z}} (g \mathrm{Z}) | f |,\tag{8}
\]

where all functions appearing in these equalities are identified with functions on G/Z.

 The subspace \(\mathcal{K}_{\chi}(\mathrm{G})\) of \(\mathcal{F}_{\chi}(\mathrm{G})\) is stable under \(\gamma_{\mathrm{G},\chi}\) and \(\delta_{\mathrm{G},\chi}\) . Let \(p\) a real number \(\geqslant 1\) . The formulas (8) imply that the representations \(\gamma_{\mathrm{G},\chi}\) and \(\delta_{\mathrm{G},\chi}\) restricted to \(\mathcal{K}_{\chi}(\mathrm{G})\) , extend to continuous, isometric linear representations of G into \(\mathcal{L}_{\chi}^{p}(\mathrm{G})\) which will be denoted \(\gamma_{\mathrm{G},\chi}^{(p)}\) and \(\delta_{\mathrm{G},\chi}^{(p)}\) . By passing to quotients, these representations also define isometric representations of G in \(\mathrm{L}_{\chi}^{p}(\mathrm{G})\) denoted in the same way.

 The representations \(\gamma_{\mathrm{G},\chi}^{(2)}\) and \(\delta_{\mathrm{G},\chi}^{(2)}\) in \(\mathrm{L}_{\chi}^{2}(\mathrm{G})\) are unitary, and will be denoted simply \(\gamma_{\mathrm{G},\chi}\) and \(\delta_{\mathrm{G},\chi}\) , respectively, when there is no confusion with the representations in \(\mathcal{F}_{\chi}(\mathrm{G})\) . We also denote \(\varrho_{\mathrm{G},\chi}\) the continuous representation of \(\mathrm{G} \times \mathrm{G}\) in \(\mathcal{L}_{\chi}^{2}(\mathrm{G})\) or \(\mathrm{L}_{\chi}^{2}(\mathrm{G})\) defined by

\[
\pmb {\varrho} _ {\mathrm{G}, \chi} (g, h) = \pmb {\gamma} _ {\mathrm{G}, \chi} (g) \circ \pmb {\delta} _ {\mathrm{G}, \chi} (h) = \pmb {\delta} _ {\mathrm{G}, \chi} (h) \circ \pmb {\gamma} _ {\mathrm{G}, \chi} (g)
\]

for all \((g,h)\in \mathrm{G}\times \mathrm{G}\) (cf. Lemma 1 of V, p. 377).

The representation \(\gamma_{\mathrm{G},\chi}\) onto \(\mathrm{L}_{\chi}^{2}(\mathrm{G})\) is none other than the induced representation \(\mathrm{Ind}_{\mathrm{Z}}^{\mathrm{G}}(\overline{\chi})\) .

When \(Z = \{e\}\) the following lemma results from INT, VIII, p. 166, § 4, no. 5, prop. 12, since G is unimodular.

Lemma 8. --- Let \(f_1 \in \mathcal{K}_\chi(\mathrm{G})\) and \(f_2 \in \mathcal{L}_\chi^2(\mathrm{G})\) .

\begin{enumerate}
\def\labelenumi{\alph{enumi})}
\item
  The function \(f\) onto \(\mathrm{G}\) defined by \(f(g) = \langle f_1 | \boldsymbol{\gamma}_{\mathrm{G},\chi}(g)f_2 \rangle\) for all \(g \in \mathrm{G}\) belongs to \(\mathcal{L}_{\chi}^{2}(\mathrm{G})\) and satisfies \(\mathrm{N}_2(f) \leqslant \mathrm{N}_1(f_1)\mathrm{N}_2(f_2)\) ;
\item
  The function \(f\) onto \(G\) defined by \(f(g) = \langle f_1 | \delta_{G,\chi}(g)f_2 \rangle\) for all \(g \in G\) belongs to \(\mathcal{L}_{\chi}^{2}(G)\) and satisfies \(N_2(f) \leqslant N_1(f_1)N_2(f_2)\) .
\end{enumerate}

Let us prove \(a\) ), the proof of assertion \(b\) ) being similar. The function \(f\) is continuous, hence \(\mu\) -measurable. For every \(z \in \mathbf{Z}\) and every \(g \in \mathbf{G}\) , we have

\[
f (g z) = \langle f _ {1} | \gamma_ {\mathrm{G}, \chi} (g z) f _ {2} \rangle = \overline {{\chi (z)}} \langle f _ {1} | \gamma_ {\mathrm{G}, \chi} (g) f _ {2} \rangle
\]

(formula (7), p. 417), hence \(f \in \mathcal{F}_{\overline{\chi}}(\mathrm{G})\) . Prop. 5 of V, p. 413 implies that it now suffices to prove that \(\mathrm{N}_2(f) \leqslant \mathrm{N}_1(f_1)\mathrm{N}_2(f_2)\) .

Suppose first that \(f_2\) belongs to \(\mathcal{K}_{\chi}(\mathrm{G})\) . For every \(g \in \mathrm{G}\) , one has by definition

\[
f (g) = \int_ {\mathrm{G/Z}} \overline {{f _ {1}}} \gamma_ {\mathrm{G}, \chi} (g) f _ {2} d \nu ,
\]

where the function \(\overline{f_1} \gamma_{\mathrm{G},\chi}(g)f_2\) is identified with a function on G/Z.

Let us define a function \(f_3\) onto \(G\) by setting \(f_3(g) = 0\) if \(f_1(g) = 0\) and \(f_3(g) = f_1(g)|f_1(g)|^{-1/2}\) otherwise. The function \(f_3\) belongs to \(\mathcal{F}_{\chi}(G)\) and satisfies \(f_1 = |f_1|^{1/2}f_3\) ; it is \(\mu\) -measurable and zero outside a compact modulo \(Z\) since \(f_1\) is. Since \(|f_1|^{1/2} \in \mathcal{K}_1(G)\) , it follows that

\[
\overline {{f _ {1}}} \gamma_ {\mathrm{G}, \chi} (g) f _ {2} = | f _ {1} | ^ {1 / 2} \overline {{f _ {3}}} \gamma_ {\mathrm{G}, \chi} (g) f _ {2},
\]

\[
| \overline {{f _ {3}}} \gamma_ {\mathrm{G}, \chi} (g) f _ {2} | = | f _ {3} | | \gamma_ {\mathrm{G}, \chi} (g) f _ {2} |
\]

for all \(g \in \mathbf{G}\) .

 Let \(g \in G\) . By the Cauchy--Schwarz inequality, we have

\[
\begin{array}{l} | f (g) | ^ {2} \leqslant \Big (\int_ {\mathrm{G/Z}} | f _ {1} | ^ {1 / 2} | f _ {3} | | \boldsymbol {\gamma} _ {\mathrm{G}, \chi} (g) f _ {2} | d \nu \Big) ^ {2} \\ \leqslant \Big (\int_ {\mathrm{G/Z}} | f _ {1} | d \nu \Big) \Big (\int_ {\mathrm{G/Z}} | f _ {3} | ^ {2} | \boldsymbol {\gamma} _ {\mathrm{G}, \chi} (g) f _ {2} | ^ {2} d \nu \Big). \end{array}
\]

Since we have \(|f_{3}|^{2}=|f_{1}|\) in \(\mathcal{K}(\mathrm{G}/\mathrm{Z})\) , we deduce by integrating over G/Z that

\[
\int_ {\mathrm{G} / \mathrm{Z}} | f (g) | ^ {2} d \nu (g) \leqslant \mathrm{N} _ {1} (f _ {1}) \int_ {\mathrm{G} / \mathrm{Z}} \left(\int_ {\mathrm{G} / \mathrm{Z}} | f _ {1} (x) | | f _ {2} (g ^ {- 1} x) | ^ {2} d \nu (x)\right) d \nu (g).
\]

The function on G/Z × G/Z deduced from the function

\[
(g, x) \mapsto | f _ {1} (x) | | f _ {2} (g ^ {- 1} x) | ^ {2}
\]

by passing to quotients is \((\nu \otimes \nu)\) measurable and compactly supported, hence it is \((\nu \otimes \nu)\) -moderate (INT, V, p. 4, § 5, no. 2, def. 2). By prop. 7 of INT, V, p. 93, § 8, no. 3, it follows that

\[
\begin{array}{l} \int_ {\mathrm{G} / \mathrm{Z}} \Bigl (\int_ {\mathrm{G} / \mathrm{Z}} | f _ {1} (x) |   | f _ {2} (g ^ {- 1} x) | ^ {2} d \nu (x) \Bigr) d \nu (g) \\ = \int_ {\mathrm{G} / \mathrm{Z}} | f _ {1} (x) | \Bigl (\int_ {\mathrm{G} / \mathrm{Z}} | f _ {2} (g ^ {- 1} x) | ^ {2} d \nu (g) \Bigr) d \nu (x) = \mathrm{N} _ {1} (f _ {1}) \mathrm{N} _ {2} (f _ {2}) ^ {2}. \end{array}
\]

Consequently, we have \(\mathrm{N}_{2}(f)^{2}\leqslant\mathrm{N}_{1}(f_{1})^{2}\mathrm{N}_{2}(f_{2})^{2}\) which establishes the desired property in this case.

 Let us consider the general case. Denote by \(u\) the linear map from \(\mathcal{K}_{\chi}(\mathrm{G})\) in \(\mathcal{L}_{\overline{\chi}}^{2}(\mathrm{G})\) which to \(f_{2}\) associates \(f\) . Let \(f_{2} \in \mathcal{L}_{\chi}^{2}(\mathrm{G})\) and let \((f_{2,n})_{n \in \mathbf{N}}\) a sequence in \(\mathcal{K}_{\chi}(\mathrm{G})\) which converges to \(f_{2}\) in \(\mathcal{L}_{\chi}^{2}(\mathrm{G})\) . Let \(f_{n} = u(f_{2,n})\) ; the sequence \((f_{n})_{n \in \mathbf{N}}\) is Cauchy in \(\mathcal{L}_{\overline{\chi}}^{2}(\mathrm{G})\) since the previous case implies that \(\mathrm{N}_{2}(f_{n} - f_{m}) \leqslant \mathrm{N}_{1}(f_{1})\mathrm{N}_{2}(f_{2,n} - f_{2,m})\) for all \(n\) and \(m\) in \(\mathbf{N}\) . Let \(f \in \mathcal{L}_{\overline{\chi}}^{2}(\mathrm{G})\) such that \((f_{n})\) converges to \(f\) (prop. 3, \(c\) ) of V, p. 409). Since \(\mathrm{N}_{2}(f_{n}) \leqslant \mathrm{N}_{1}(f_{1})\mathrm{N}_{2}(f_{2,n})\) for all \(n \in \mathbf{N}\) , it follows that \(\mathrm{N}_{2}(f) \leqslant \mathrm{N}_{1}(f_{1})\mathrm{N}_{2}(f_{2})\) .

There exists a subsequence \((f_{n_{k}})_{k\in\mathbf{N}}\) such that \(f_{n_{k}}(g)\) converges to \(f(g)\) for all \(g\in G\) outside of a negligible subset of G modulo Z (loc. cit., b)). But on the other hand, for every \(g\in G\) , we have

\[
f _ {n _ {k}} (g) = \left\langle f _ {1} \mid \gamma_ {\mathrm{G}, \chi} (g) f _ {2, n _ {k}} \right\rangle\rightarrow \left\langle f _ {1} \mid \gamma_ {\mathrm{G}, \chi} (g) f _ {2} \right\rangle .
\]

Consequently, we have \(f(g) = \langle f_1 | \gamma_{\mathrm{G},\chi}(g)f_2 \rangle\) for all \(g\) outside a negligible subset of G modulo Z. Since \(f \in \mathcal{L}_{\overline{\chi}}^2(\mathrm{G})\) and that \(\mathrm{N}_2(f) \leqslant \mathrm{N}_1(f_1)\mathrm{N}_2(f_2)\) , the lemma is proved.

